% ============================================================================
%  preamble.tex — toàn bộ cấu hình định dạng của luận văn
%  Biên dịch: XeLaTeX (Overleaf: Menu → Compiler → XeLaTeX)
%  Thư mục tham chiếu: biblatex + biber (Overleaf tự nhận diện)
%
%  Nếu trường có template riêng, chỉ cần thay THẾ file này — nội dung các
%  chương trong chapters/ không phụ thuộc vào định dạng.
% ============================================================================

% ---------------------------------------------------------------------------
% 1. Trang giấy và lề
% ---------------------------------------------------------------------------
% Cỡ chữ 13pt — đặt qua gói `fontsize` vì lớp report chuẩn chỉ nhận 10/11/12pt
% còn extsizes thì thiếu đúng mức 13pt. Gói này cho phép cỡ chữ tuỳ ý và điều
% chỉnh đồng bộ toàn bộ các cỡ dẫn xuất (\large, \small, chú thích, chỉ số trên...).
\usepackage[fontsize=13pt]{fontsize}

% Cho phép các font Type1 (Computer Modern dùng trong công thức toán) nhận cỡ
% chữ tuỳ ý. Thiếu gói này, mọi cỡ không nằm trong bảng dựng sẵn — trong đó có
% đúng 13pt — đều bị thay thế kèm cảnh báo "Font shape ... not available".
\usepackage{anyfontsize}

\usepackage{geometry}
\geometry{
    a4paper,
    left=3.5cm,      % lề trái rộng để đóng bìa
    right=2.0cm,
    top=3.0cm,
    bottom=3.0cm,
    footskip=1.2cm
}

% ---------------------------------------------------------------------------
% 2. Font và ngôn ngữ tiếng Việt (XeLaTeX)
% ---------------------------------------------------------------------------
\usepackage{fontspec}

% Font được nạp theo TÊN TỆP (không theo tên hệ thống) để chạy được trên mọi
% môi trường: Overleaf, TeX Live, MiKTeX và Tectonic. Tectonic không dùng
% fontconfig nên tra theo tên hệ thống ("TeX Gyre Termes") sẽ thất bại.
% Ba họ font dưới đây là bản tự do tương đương Times, Helvetica và Courier.
\setmainfont{texgyretermes}[
    Extension      = .otf,
    UprightFont    = *-regular,
    BoldFont       = *-bold,
    ItalicFont     = *-italic,
    BoldItalicFont = *-bolditalic,
]
\setsansfont{texgyreheros}[
    Extension      = .otf,
    UprightFont    = *-regular,
    BoldFont       = *-bold,
    ItalicFont     = *-italic,
    BoldItalicFont = *-bolditalic,
]
\setmonofont{texgyrecursor}[
    Extension      = .otf,
    UprightFont    = *-regular,
    BoldFont       = *-bold,
    ItalicFont     = *-italic,
    BoldItalicFont = *-bolditalic,
    Scale          = 0.92,
]

\usepackage{polyglossia}
\setdefaultlanguage{vietnamese}
\setotherlanguage{english}

% Lưới an toàn: ép lại tên tiếng Việt phòng khi gói ngôn ngữ dịch thiếu.
% Dùng \AtBeginDocument thay cho \renewcaptionname để tương thích với MỌI
% phiên bản polyglossia — lệnh \renewcaptionname chỉ có ở các bản mới.
\AtBeginDocument{%
    \renewcommand{\chaptername}{Chương}%
    \renewcommand{\figurename}{Hình}%
    \renewcommand{\tablename}{Bảng}%
    \renewcommand{\contentsname}{Mục lục}%
    \renewcommand{\listfigurename}{Danh mục hình vẽ}%
    \renewcommand{\listtablename}{Danh mục bảng biểu}%
    \renewcommand{\appendixname}{Phụ lục}%
    \renewcommand{\indexname}{Chỉ mục}%
}

% ---------------------------------------------------------------------------
% 3. Giãn dòng và đoạn văn
% ---------------------------------------------------------------------------
\usepackage{setspace}
\onehalfspacing                        % giãn dòng 1,5
\setlength{\parindent}{1cm}
\setlength{\parskip}{0.35em}

% Nới ràng buộc dàn dòng. Văn bản tiếng Việt có nhiều từ đơn âm tiết ngắn và
% hầu như không ngắt từ được, lại xen kẽ nhiều thuật ngữ tiếng Anh cùng tên
% hàm, tên tệp dài — nếu giữ mặc định sẽ còn khá nhiều dòng tràn ra lề.
% \emergencystretch cho TeX một lượng giãn dự phòng ở lượt dàn dòng cuối cùng.
\tolerance=2000
\setlength{\emergencystretch}{3em}
\hbadness=6000                         % bớt cảnh báo dòng thưa chữ không đáng kể

% ---------------------------------------------------------------------------
% 4. Toán học
% ---------------------------------------------------------------------------
\usepackage{amsmath}
\usepackage{amssymb}
\usepackage{amsfonts}

% ---------------------------------------------------------------------------
% 5. Hình ảnh và bảng biểu
% ---------------------------------------------------------------------------
\usepackage{graphicx}
\graphicspath{{figures/}}
\usepackage{float}
\usepackage{booktabs}
\usepackage{multirow}
\usepackage{longtable}
\usepackage{array}
\usepackage{tabularx}
\usepackage{caption}
\usepackage{subcaption}
\captionsetup{font=small, labelfont=bf, skip=6pt}
\captionsetup[table]{position=above}

% Cột canh giữa/phải có độ rộng cố định
\newcolumntype{C}[1]{>{\centering\arraybackslash}p{#1}}
\newcolumntype{R}[1]{>{\raggedleft\arraybackslash}p{#1}}
\newcolumntype{L}[1]{>{\raggedright\arraybackslash}p{#1}}

% ---------------------------------------------------------------------------
% 6. Màu và sơ đồ TikZ (tái sử dụng từ paper_ijai.tex)
% ---------------------------------------------------------------------------
\usepackage{xcolor}
\usepackage{tikz}
\usetikzlibrary{arrows.meta, positioning, fit, calc}

\definecolor{stageblue}{HTML}{DCEBFF}
\definecolor{stagegreen}{HTML}{DDF4E7}
\definecolor{stageorange}{HTML}{FFE8CC}
\definecolor{stagegray}{HTML}{EEF1F5}
\definecolor{stagepurple}{HTML}{EDE4FF}
\definecolor{arrowgray}{HTML}{4B5563}

\tikzset{
    pipelinebox/.style={
        draw=arrowgray, rounded corners=2pt, align=center,
        minimum width=3.1cm, minimum height=1.15cm,
        font=\small, inner sep=5pt
    },
    data/.style={pipelinebox, fill=stagegray},
    classmodel/.style={pipelinebox, fill=stageblue},
    xai/.style={pipelinebox, fill=stageorange},
    sam/.style={pipelinebox, fill=stagepurple},
    segmodel/.style={pipelinebox, fill=stagegreen},
    flow/.style={-{Latex[length=2.6mm]}, thick, draw=arrowgray},
    dashedflow/.style={-{Latex[length=2.6mm]}, thick, dashed, draw=arrowgray},
    lane/.style={draw=arrowgray!45, rounded corners=4pt, inner sep=8pt}
}

% Bảng màu Okabe–Ito, đồng bộ với màu trong toàn bộ notebook
\definecolor{okOrange}{HTML}{E69F00}
\definecolor{okSky}{HTML}{56B4E9}
\definecolor{okGreen}{HTML}{009E73}
\definecolor{okBlue}{HTML}{0072B2}
\definecolor{okVermil}{HTML}{D55E00}

% ---------------------------------------------------------------------------
% 7. Đánh số chương mục và tiêu đề
% ---------------------------------------------------------------------------
\usepackage{titlesec}
\titleformat{\chapter}[display]
    {\normalfont\Large\bfseries\centering}
    {\chaptertitlename\ \thechapter}{12pt}{\LARGE}
\titlespacing*{\chapter}{0pt}{-10pt}{28pt}

\setcounter{secnumdepth}{3}
\setcounter{tocdepth}{2}

% ---------------------------------------------------------------------------
% 8. Header / footer
% ---------------------------------------------------------------------------
\usepackage{fancyhdr}
\pagestyle{fancy}
\fancyhf{}
\fancyfoot[C]{\thepage}
\renewcommand{\headrulewidth}{0pt}
\fancypagestyle{plain}{\fancyhf{}\fancyfoot[C]{\thepage}\renewcommand{\headrulewidth}{0pt}}

% ---------------------------------------------------------------------------
% 9. Liên kết chéo và trích dẫn
% ---------------------------------------------------------------------------
% Cho phép ngắt dòng bên trong văn bản đánh máy (\texttt). Mặc định LaTeX đặt
% \hyphenchar = -1 cho font đánh máy nên các chuỗi dài như tên tệp, đường dẫn
% hay định danh mã nguồn không thể ngắt và tràn ra ngoài lề.
\usepackage[htt]{hyphenat}

\usepackage[hidelinks, unicode]{hyperref}

% xurl cho phép ngắt URL tại bất kỳ ký tự nào. Cần thiết vì một số mục trong
% danh mục tham khảo chứa địa chỉ kho mã nguồn dài. Phải nạp sau hyperref.
\usepackage{xurl}
\hypersetup{
    pdftitle={Phat hien benh la sau rieng bang hoc sau ket hop XAI va pseudo-labeling},
    pdfauthor={},
    bookmarksnumbered=true
}
% Tuỳ chọn `english` ép cleveref nạp bộ định nghĩa tiếng Anh thay vì tự dò theo
% babel/polyglossia. Cần thiết vì cleveref 0.21.x không biết ngôn ngữ
% `vietnamese`; nếu để nó tự dò thì các macro liên từ (\crefpairconjunction...)
% sẽ không được định nghĩa và mọi lệnh \cref nhiều nhãn đều lỗi.
% Toàn bộ chuỗi hiển thị được Việt hoá ngay bên dưới nên tuỳ chọn này không
% làm lộ chữ tiếng Anh nào ra bản in.
\usepackage[nameinlink, english]{cleveref}

% Bản vá tương thích cleveref 0.21.x với polyglossia 1.5x: cleveref chỉ đặt
% \cref@language khi \languagename đã tồn tại, còn polyglossia bản mới định
% nghĩa \languagename muộn hơn. Hệ quả là \cref@language chưa xác định khi
% cleveref chạy hook \AtBeginDocument, gây lỗi "Undefined control sequence".
% Khai báo dự phòng dưới đây vô hại nếu cleveref đã tự đặt được giá trị.
\makeatletter
\providecommand{\cref@language}{english}
\makeatother

% Việt hoá các liên từ khi tham chiếu nhiều nhãn cùng lúc:
%   \cref{a,b}   -> "Hình 1 và Hình 2"
%   \cref{a,b,c} -> "Hình 1, Hình 2 và Hình 3"
%   \crefrange   -> "Hình 1 đến Hình 5"
%
% Phải đặt trong \AtBeginDocument: cleveref chỉ tạo các macro liên từ ở hook
% \AtBeginDocument của chính nó (sau khi đã xác định ngôn ngữ), nên ở thời điểm
% preamble chúng chưa tồn tại và \renewcommand sẽ báo "Command undefined".
% Hook khai báo tại đây chạy sau hook của cleveref vì được đăng ký muộn hơn.
\newcommand{\vnand}{ và~}
\AtBeginDocument{%
    \renewcommand{\crefpairconjunction}{\vnand}%
    \renewcommand{\crefmiddleconjunction}{, }%
    \renewcommand{\creflastconjunction}{\vnand}%
    \renewcommand{\crefpairgroupconjunction}{\vnand}%
    \renewcommand{\crefmiddlegroupconjunction}{, }%
    \renewcommand{\creflastgroupconjunction}{\vnand}%
    \renewcommand{\crefrangeconjunction}{ đến~}%
    \renewcommand{\crefrangepreconjunction}{}%
    \renewcommand{\crefrangepostconjunction}{}%
}

\crefname{figure}{Hình}{Hình}
\crefname{table}{Bảng}{Bảng}
\crefname{chapter}{Chương}{Chương}
\crefname{section}{Mục}{Mục}
\crefname{equation}{Công thức}{Công thức}

% Dùng backend=bibtex thay vì biber: bibtex có sẵn trong Tectonic và mọi bản
% TeX Live/MiKTeX, trong khi biber là chương trình ngoài phải cài riêng.
% Tệp refs.bib chỉ chứa ký tự ASCII và các dấu phụ đã escape theo cú pháp
% LaTeX (\'e, \"a...) nên bibtex xử lý được đầy đủ.
% Khai báo language=english để biblatex dùng bộ chuỗi tiếng Anh cho phần tài
% liệu tham khảo — phù hợp vì toàn bộ tài liệu trích dẫn đều là công bố quốc tế.
% autolang=none: chặn biblatex tự đổi ngôn ngữ theo ngôn ngữ tài liệu. Không có
% tuỳ chọn này, biblatex sẽ cố dựng mục tài liệu tham khảo bằng tiếng Việt,
% nhưng không tồn tại tệp vietnamese.lbx nên hàng loạt chuỗi ("and", "pages",
% "in"...) bị bỏ chưa dịch và hiện ra thô trong bản in.
\usepackage[backend=bibtex, style=ieee, sorting=none,
            language=english, autolang=none]{biblatex}

% biblatex chốt ngôn ngữ mục tài liệu tham khảo ngay lúc nạp gói, dựa theo ngôn
% ngữ chính của tài liệu. Vì không tồn tại tệp vietnamese.lbx, cần ánh xạ tường
% minh sang english.lbx — đây là cơ chế chính thức của biblatex cho các ngôn ngữ
% chưa có bộ chuỗi riêng. Thiếu dòng này, các chuỗi "and", "pages", "in",
% "jourvol"... sẽ bị bỏ chưa dịch và hiện thô trong danh mục tham khảo.
\DeclareLanguageMapping{vietnamese}{english}

\addbibresource{refs.bib}

% ---------------------------------------------------------------------------
% 10. Ghi chú việc còn thiếu — mọi chỗ số liệu chưa xác minh phải dùng \todo
% ---------------------------------------------------------------------------
% \TODO được cố tình định nghĩa ở dạng NỘI DÒNG (không dùng todonotes) để có thể
% đặt ở bất kỳ đâu, kể cả giữa các dòng của trang bìa — nơi lệnh \todo[inline]
% của todonotes sẽ phá vỡ cấu trúc đoạn và gây lỗi "There's no line here to end".
% Nền vàng viền đỏ bảo đảm không thể bỏ sót khi rà soát bản in.
% KHÔNG bọc trong \colorbox: hộp màu là một hộp ngang không thể ngắt dòng, nên
% một ghi chú dài sẽ tràn ra ngoài lề hàng chục xăng-ti-mét. Chữ đỏ đậm ngắt
% dòng bình thường mà vẫn đủ nổi bật để không bỏ sót khi rà bản in.
\newcommand{\TODO}[1]{%
    \textcolor{red}{\textbf{[\,CẦN BỔ SUNG: #1\,]}}%
}

% ---------------------------------------------------------------------------
% 11. Macro dùng chung
% ---------------------------------------------------------------------------
\usepackage{siunitx}
\sisetup{output-decimal-marker={,}, group-separator={.}, group-minimum-digits=4}

% Tên kỹ thuật viết nhất quán trong toàn luận văn
\newcommand{\EffNet}{EfficientNet-B0}
\newcommand{\GradCAM}{Grad-CAM}
\newcommand{\GradCAMpp}{Grad-CAM++}
\newcommand{\UNetpp}{UNet++}
\newcommand{\SAM}{SAM}

% Thuật ngữ tiếng Anh giữ nguyên, in nghiêng ở lần xuất hiện đầu
\newcommand{\en}[1]{\textit{#1}}

\usepackage{enumitem}
\setlist{itemsep=2pt, topsep=4pt, parsep=0pt}

\usepackage{microtype}
